\documentclass[11pt]{article}
\usepackage[margin=1in]{geometry}
\usepackage{amsmath,amssymb,amsthm,mathtools,microtype,hyperref}
\hypersetup{hidelinks}
\newtheorem{theorem}{Theorem}
\newcommand{\ddc}{dd^c}
\newcommand{\C}{\mathbb C}
\newcommand{\R}{\mathbb R}
\newcommand{\Z}{\mathbb Z}
\newcommand{\PSH}{\operatorname{PSH}}
\newcommand{\vol}{\operatorname{vol}}
\newcommand{\NP}[1]{\langle #1\rangle}
\title{Partial kernels from Newton and multiplier-ideal series}
\author{}
\date{4 October 2026}
\begin{document}
\maketitle
\begin{abstract}
We prove sufficient conditions for the envelope lower bound and normalized
kernel-measure convergence of a specified graded section series. On a
product of projective lines we select precisely the monomials in an interior
rational Newton polytope, sample over the whole manifold with a smooth
positive invariant probability, and use a smooth invariant obstacle whose
logarithmic weight is strictly convex. We prove the entire global selection
constraint, identify its envelope among all genuine plurisubharmonic
competitors, derive Bernstein--Markov, and compute the growth dimension
from the actual spaces. The limit is relative lattice volume lifted by the
inverse weight gradient and angular Haar measure. Lower-dimensional
selection gives a positive singular measure although the top
nonpluripolar curvature power vanishes. An explicit rational-map example
admits a nonclosed sampling background. Independently, an actual
multiplier-ideal series on a product of two projective lines has
Hilbert growth exponent one and three adjacent monomial bands. We
identify its full ideal envelope among arbitrary competitors, prove
equality with the ordinary upper type class, and compute its original-norm
kernel limits. Its intrinsic limiting mass is three; pole-sensitive
bounds give global lower convergence and uniform convergence after
exponentiation. The general selecting-series programme remains open.
\end{abstract}

\section{The specified series and its target}
This paper answers the sufficient-condition part of catalogue problem P20.
Use $\ddc=(i/(2\pi))\partial\bar\partial$. A $\theta$-psh function
means a genuine quasi-plurisubharmonic function with its canonical upper
representative: in every holomorphic frame, $v+\varphi_h$ is
plurisubharmonic and satisfies the submean inequality. Positive
distributional curvature alone does not authorize pointwise changes.

Let $X=(\mathbb P^1)^n$, $n\geq1$, and
$L=\mathcal O(k_1,\ldots,k_n)$, with positive integers $k_j$ and product
Fubini--Study metric $h$. On $T=(\C^*)^n$, put
\[
 r_j=\log|z_j|^2,\qquad
 \varphi_0(r)=\sum_jk_j\log(1+e^{r_j}),\qquad
 \theta=c_1(L,h)=\ddc\varphi_0.
\]
The ambient Hermitian form $\omega$ may be any smooth positive form;
it is not substituted for the closed Chern curvature $\theta$.
Write $P=\prod_j[0,k_j]$. Fix a nonempty rational polytope
$Q\Subset\operatorname{int}P$, with
\[
 \operatorname{aff}Q=b+W,\qquad b\in\Z^n,\qquad
 W\text{ a rational vector space},\qquad d=\dim Q.
\]
Set $\Lambda=W\cap\Z^n$. Relative volume $dq_\Lambda$ gives a
fundamental cell of $\Lambda$ volume one; a singleton has volume one
when $d=0$. The selected spaces are exactly
\begin{equation}\label{eq:spaces}
 V_m=\operatorname{span}\{\sigma_\alpha=z^\alpha e^m:
                   \alpha\in mQ\cap\Z^n\}.
\end{equation}
These sections extend globally because $0\leq\alpha_j\leq mk_j$.
Let $w$ be smooth, real and torus-invariant on $X$, and assume
\begin{equation}\label{eq:convex}
 F=\varphi_0+w,\qquad D^2F(r)>0\quad(r\in\R^n).
\end{equation}
There is no uniform Hessian bound at infinity in this assumption.
The choice $w=0$ satisfies it. Let $\nu$ be a smooth strictly positive
invariant probability on the whole of $X$. On $T$,
\[
 d\nu=a(r)\,dr\,d\tau,\qquad a>0,\quad
 \int_{\R^n}a\,dr=1,
\]
where $d\tau$ is angular Haar probability. This includes
$\nu=\omega^n/\int_X\omega^n$ when $\omega$ is invariant, even if
it is nonclosed. For a general $\omega$ an invariant sampling
probability can be chosen independently.

The Hilbert norm, kernel and logarithmic potential are
\[
 \|s\|_m^2=\int_X|s|_{h^m}^2e^{-mw}\,d\nu,\qquad
 B_m=\sum_j|s_{m,j}|_{h^m}^2,\qquad U_m=m^{-1}\log B_m,
\]
for an orthonormal basis of the actual $V_m$. Put
$H_Q(r)=\max_{q\in Q}q\cdot r$. The constraint and target are
\begin{align}
 \mathcal C_Q&=\{v\in\PSH(X,\theta):
        v+\varphi_0\leq H_Q+C_v\text{ on }T,\ C_v<\infty\},\notag\\
 E_Q&=\left(\sup\{v\in\mathcal C_Q:
                 v\leq w\text{ quasi-everywhere on }X\}\right)^*.
                 \label{eq:envelope}
\end{align}
Functions in $\mathcal C_Q$ are not identically $-\infty$;
quasi-everywhere excludes a pluripolar set. This is a Newton-growth
constraint; identification with the multiplier-ideal envelope is not
assumed.

\begin{theorem}\label{thm:main}
Under these hypotheses $\nabla F:\R^n\to\operatorname{int}P$ is a
smooth diffeomorphism. Put $r(q)=(\nabla F)^{-1}(q)$,
$F^*(q)=q\cdot r(q)-F(r(q))$ and
\[
 G_Q(r)=\max_{q\in Q}\{q\cdot r-F^*(q)\}.
\]
Then $E_Q=G_Q-\varphi_0$ on $T$, with canonical value $-\infty$
on $X\setminus T$. The norms are positive definite and satisfy weighted
Bernstein--Markov. The actual dimensions $N_m=\dim V_m$ satisfy
\[
 N_m=m^d\vol_\Lambda(Q)+O(m^{d-1})\quad(d\geq1),
 \qquad N_m=1\quad(d=0).
\]
For all sufficiently large $m$,
\begin{equation}\label{eq:error}
 |U_m-E_Q|\leq\frac{C}{m}
                \left(1+\sum_j|r_j|+\log m\right)\quad\text{on }T.
\end{equation}
Thus convergence is locally uniform on $T$, in every finite $L^p$
against smooth volume on $X$, and weakly for closed curvature currents.
Both potentials are $-\infty$ on $X\setminus T$, so the extended
pointwise lower bound holds on all $X$.

The normalized finite Borel measures converge against every continuous
test function on $X$:
\begin{equation}\label{eq:limit}
 \mu_m=d!m^{-d}e^{-mw}B_m\nu\ \rightharpoonup\
 \mu_Q=\Psi_*(d!dq_\Lambda\,d\tau),\qquad
 \Psi(q,\tau)=e^{r(q)/2}\tau.
\end{equation}
Their limiting mass is $d!\vol_\Lambda(Q)>0$. For $d=0$ this is
Haar probability on one real $n$-torus. For $d<n$ the limit is singular
with respect to every smooth ambient volume and
$\NP{(\theta+\ddc E_Q)^{d+1}}=0$. For $d=n$,
\begin{equation}\label{eq:fullma}
 \mu_Q=\NP{(\theta+\ddc E_Q)^n}.
\end{equation}
\end{theorem}

\section{Global production and all envelope competitors}
Every section of $L^m$ is a polynomial in the standard affine frame,
with degree at most $mk_j$ in $z_j$, as follows by changing to inverted
coordinates at infinity. More precisely,
\begin{equation}\label{eq:entire}
 s\in V_m\quad\Longleftrightarrow\quad
 |s/e^m|^2\leq C_s e^{mH_Q(r)}\text{ on }T
 \quad\text{for some finite }C_s.
\end{equation}
The forward implication is the triangle inequality. For the converse,
write $s/e^m=\sum_\alpha c_\alpha z^\alpha$ and average over angles:
\[
 \sum_\alpha|c_\alpha|^2e^{\alpha\cdot r}
                   \leq C_s e^{mH_Q(r)}.
\]
If a nonzero coefficient had $\alpha\notin mQ$, separation gives
$u$ with $\alpha\cdot u>mH_Q(u)$; substituting $r=tu$ and letting
$t\to\infty$ contradicts this inequality. This proves the entire
constraint, and convexity gives $V_mV_l\subset V_{m+l}$.

For $p\in\operatorname{int}P$, boundedness of $w$ gives
\begin{equation}\label{eq:coercive}
 F(r)-p\cdot r\geq\delta_p\sum_j|r_j|-\|w\|_\infty,
 \qquad\delta_p=\min_j(p_j,k_j-p_j)>0.
\end{equation}
Strict convexity and coercivity give a unique minimizer, with
$\nabla F(r)=p$. Conversely a supporting plane and the asymptotic
slopes of $\varphi_0$ imply that every gradient belongs to $P$.
The image is open by the inverse function theorem, so belongs to
$\operatorname{int}P$. This proves the diffeomorphism.
Moreover $\nabla F^*=r$ and $D^2F^*=(D^2F(r(q)))^{-1}$.
These derivatives are bounded on a neighborhood of $Q$, and $r(Q)$
is compact.

In a product chart with some coordinates inverted, the candidate's
local potential is
\begin{equation}\label{eq:charts}
 \max_{q\in Q}\left\{\sum_j c_j(q)\log|\zeta_j|^2-F^*(q)\right\},
 \qquad c_j(q)=q_j\text{ or }k_j-q_j.
\end{equation}
All coefficients are bounded below by a fixed $\delta>0$. The supremum
is locally bounded above, continuous away from coordinate hyperplanes,
and tends uniformly to $-\infty$ at each hyperplane with other
coordinates bounded. It is locally integrable by comparison with a
fixed logarithmic summand. It is genuinely psh: at a finite center use
a maximizing summand and its submean inequality; at a $-\infty$ center
the inequality is automatic. These potentials glue by frame changes.
Boundedness of $F^*$ gives $G_Q-H_Q=O(1)$, and Fenchel's inequality
gives $G_Q\leq F$. Thus the candidate belongs to the envelope class and
is at most $w$ everywhere.

Take any competitor $v$, without assuming invariance. Its
quasi-everywhere obstacle is an everywhere obstacle: pluripolar sets
are locally Lebesgue-null, and shrinking-ball submeans pass
$v+\varphi_h\leq w+\varphi_h$ from almost everywhere to every point.
On $T$ form the circular maximum
\[
 M(r)=\max_{\tau\in(S^1)^n}\{v(e^{r/2}\tau)+\varphi_0(r)\}.
\]
Compactness makes it upper semicontinuous. Rotate a maximizing center
and apply its psh submean inequality to see that the invariant lift is
psh. This lift extends genuinely across the standard affine $\C^n$
after adding $\varphi_0$. If $M$ were $-\infty$ at some $r$,
successive one-variable maximum principles would force that psh
function to be $-\infty$ throughout the polydisc with radii $e^{r_j/2}$,
contradicting nontriviality and local integrability. Thus $M$ is finite
everywhere. In logarithmic complex coordinates it depends only on real
parts; restriction to lines or convolution gives a positive
distributional real Hessian, hence convexity.

We have $M\leq F$ and $M\leq H_Q+C_v$. Every supporting slope $p$
of $M$ lies in $Q$, by separation and growth along a ray. Such a slope
exists at every point of a finite convex function. At a given $r$,
\[
 M(r)=p\cdot r-M^*(p)\leq p\cdot r-F^*(p)\leq G_Q(r).
\]
Here $M\leq F$ implies $M^*\geq F^*$. In a boundary chart the
domination holds almost everywhere; shrinking-ball submeans and the
uniform decay in \eqref{eq:charts} give value $-\infty$ for $v$ at
the boundary. The candidate already has its canonical representative,
so the supremum and its upper regularization both equal it.

\section{Exact norms, lattice grids and the lower bound}
Angular Fourier orthogonality gives
\begin{align}
 A_m(\alpha)=\|\sigma_\alpha\|_m^2
     &=\int_{\R^n}e^{\alpha\cdot r-mF(r)}a(r)\,dr,\label{eq:norm}\\
 \Big\|\sum c_\alpha\sigma_\alpha\Big\|_m^2
     &=\sum|c_\alpha|^2A_m(\alpha).\notag
\end{align}
Every $A_m$ is finite and strictly positive. The weighted norm is also
exactly the unweighted norm for $he^{-w}$, with local potential $F$
and closed curvature. The orthonormal monomials give
\begin{equation}\label{eq:kernel}
 B_m(r)=\sum_{\alpha\in mQ\cap\Z^n}
          \frac{e^{\alpha\cdot r-m\varphi_0(r)}}{A_m(\alpha)}.
\end{equation}
For nonempty $V_m$, the local expression
$U_m+\varphi_h=m^{-1}\log\sum_\alpha|f_\alpha|^2$
is genuinely psh, so its curvature current is closed and positive.
For $q=\alpha/m$, $F(r)-q\cdot r$ has its unique minimum $-F^*(q)$
at $r(q)$. On a fixed compact neighborhood of $r(Q)$, $a$ has a
positive lower bound and $D^2F$ an upper bound. Taylor's formula
on a ball of radius $m^{-1/2}$ gives the uniform estimates
\begin{equation}\label{eq:laplace}
 c m^{-n/2}e^{mF^*(q)}\leq A_m(\alpha)\leq e^{mF^*(q)}.
\end{equation}
The upper bound uses only the minimum and $\int a=1$. No matching
upper Gaussian asymptotic is needed.

The integral offset identifies the lattice in every degree:
$(m\operatorname{aff}Q)\cap\Z^n=mb+\Lambda$. Fundamental cells
anchored in $m(Q-b)$ give the stated dimension asymptotics.
Discrepancies are confined to a fixed-width neighborhood of the finitely
many facets, with volume and cell count $O(m^{d-1})$. The same
argument proves lattice Riemann sums for continuous functions on $Q$.
For $d=0$, $Q=\{b\}$ and $N_m=1$.

There is an all-degrees grid approximation. Choose a relative interior
$q_0$ with inscribed relative ball radius $\rho>0$, and a lattice
covering radius $R_\Lambda>0$. Moving any $q\in Q$ to
$(1-t)q+tq_0$, with $t=2R_\Lambda/(m\rho)$, puts a relative ball
of radius $2R_\Lambda/m$ around it inside $Q$. Round to
$b+m^{-1}\Lambda$ to obtain a point in
$Q_m=Q\cap(b+m^{-1}\Lambda)$ at distance $O(1/m)$ from $q$.
This works for every sufficiently large degree. The singleton case
is exact.

Put $G_{Q,m}(r)=\max_{q\in Q_m}(q\cdot r-F^*(q))$.
The grid approximation and bounded gradient of $F^*$ on $Q$ give
\[
 0\leq G_Q-G_{Q,m}\leq C(1+\textstyle\sum_j|r_j|)/m.
\]
Equations \eqref{eq:kernel}--\eqref{eq:laplace} imply
\begin{equation}\label{eq:sandwich}
 G_{Q,m}-\varphi_0\leq U_m\leq G_{Q,m}-\varphi_0+
               \frac{(n/2)\log m+\log N_m-\log c}{m}.
\end{equation}
The left inequality uses one explicitly produced global monomial
approximating a maximizing slope. It respects the entire constraint
\eqref{eq:entire}. The dimension estimate proves \eqref{eq:error}.
All selected monomials vanish at the toric boundary because $Q$ is
separated from every face of the degree box. Logarithmic coordinate
singularities have finite $L^p$ norm for every finite $p$ against
smooth volume. This proves global $L^p$ convergence, and distributional
differentiation of $L^1$ convergence gives weak curvature convergence.
No uniform error across the logarithmic poles is asserted.

Fenchel's inequality and \eqref{eq:laplace} yield
$\sup_X e^{-mw}B_m\leq C N_m m^{n/2}$. Cauchy--Schwarz in the exact
orthonormal basis therefore gives
\[
 \sup_X|s|_{h^m}e^{-mw/2}
             \leq C m^{d/2+n/4}\|s\|_m\quad(s\in V_m).
\]
Finitely many initial degrees can be absorbed in $C$, including empty
spaces. A polynomial is bounded by $C_\varepsilon e^{\varepsilon m}$,
which proves weighted Bernstein--Markov.

\section{Uniform concentration and the kernel measure}
Each monomial defines a probability
\[
 dP_{m,q}=\frac{e^{m(q\cdot r-F(r))}a(r)\,dr\,d\tau}
                         {A_m(\alpha)},\qquad q=\alpha/m.
\]
For every $\eta>0$, uniqueness of minima, compactness and
\eqref{eq:coercive} give $c_\eta>0$ with
\[
 F(r)-q\cdot r+F^*(q)\geq c_\eta
       \quad\text{if }|r-r(q)|\geq\eta,\quad q\in Q.
\]
Inside a large compact set use continuity and uniqueness; outside it
use the uniform positive coercivity constant given by
$Q\Subset\operatorname{int}P$. Equation \eqref{eq:laplace} implies
\[
 P_{m,q}\{|r-r(q)|\geq\eta\}\leq C m^{n/2}e^{-mc_\eta}.
\]
For $f\in C(X)$ write $\bar f(r)=\int f(e^{r/2}\tau)d\tau$.
Uniform continuity near $r(Q)$ gives
$\int f\,dP_{m,q}\to\bar f(r(q))$ uniformly on the selected grid.
The exact identity
\[
 \mu_m=d!m^{-d}\sum_{\alpha\in mQ\cap\Z^n}P_{m,\alpha/m}
\]
and lattice Riemann sums prove \eqref{eq:limit}. In particular
$\mu_m(X)=d!m^{-d}N_m\to d!\vol_\Lambda(Q)$.
Using $n!m^{-n}$ when $d<n$ would give zero. The limit is independent
of the particular smooth positive invariant sampling probability,
although all finite-degree Hilbert norms depend on it.

\section{Ambient Hermitian geometry and curvature powers}
The map $\Psi$ embeds
$S_Q=\Psi(\operatorname{relint}Q\times(S^1)^n)$ as a real
$(n+d)$-dimensional submanifold: the inverse gradient restricted to
$b+W$ is an embedding, and the angular directions are independent.
For $d=0$ use the single torus. Give $X$ the Riemannian metric
$g_\omega(u,v)=\omega(u,Jv)$, with
$d\vol_\omega=\omega^n/n!$. Define a smooth positive Jacobian by
\[
 \Psi^*d\vol_{\omega,S_Q}=J_\omega(q,\tau)\,dq_\Lambda d\tau.
\]
The relative boundary of $Q$ is null, so the geometric density is
\begin{equation}\label{eq:jacobian}
 d\mu_Q=\frac{d!}{J_\omega\circ\Psi^{-1}}\,d\vol_{\omega,S_Q}.
\end{equation}
No closedness of $\omega$ is needed. For $d<n$, the closure of
$S_Q$ has ambient volume zero and carries the entire positive limit.

The restriction of $F^*$ to $Q$ is uniformly strongly convex in $W$.
The maximizing $q(r)$ is unique, and variational inequalities give
\[
 (r-s)\cdot(q(r)-q(s))
 \geq(\nabla F^*(q(r))-\nabla F^*(q(s)))\cdot(q(r)-q(s))
 \geq c|q(r)-q(s)|^2.
\]
Hence $q$ is Lipschitz. Comparing maximizing values at $r$ and $r+h$
gives $\nabla G_Q=q$, so $G_Q$ is $C^{1,1}$. Its almost-everywhere
Hessian has rank at most $d$, because its gradient lies in $b+W$.
The differentiability and weak-derivative identification used here follow
componentwise from \cite[Theorems 3.28, 3.31 and Corollary 3.32]{Kinnunen};
local cutoffs suffice on each logarithmic chart.

On $T$, $T_Q=\theta+\ddc E_Q=\ddc G_Q$. Mollify $G_Q$ in
logarithmic coordinates on compact subsets. Convex mollifications
converge in $C^1$, have uniformly bounded Hessians, and their Hessians
converge almost everywhere and in every finite local $L^p$.
Every Hessian minor therefore converges in $L^1$. The bounded-psh
wedge agrees with these minors: induct on the definition
$\ddc(G_Q T_Q^{j-1})$, using uniform convergence, integration by
parts and local bounds on wedge coefficients. The nonpluripolar
extension agrees with this product on $T$ and gives no mass to the
pluripolar boundary, by \cite[Section 1.2, Proposition 1.6 and
Theorem 1.8]{BEGZ}. Their hypotheses hold: $X$ is compact K\"ahler,
$\theta$ is K\"ahler, $E_Q$ is genuine $\theta$-psh, and its unbounded
locus is the analytic toric boundary. The comparison of smooth positive
forms on compact $X$ also makes all Hermitian trace measures finite.
Thus $\NP{T_Q^{d+1}}=0$ for $d<n$.
Also $\NP{T_Q^d}\wedge\omega^{n-d}$ has locally bounded densities
on $T$, so gives $S_Q$ zero mass. It cannot equal the positive singular
kernel limit for any smooth ambient $\omega$; when $d=0$ its zeroth
power is interpreted as $1$.

For $d=n$, on $r(\operatorname{int}Q)$ the maximizer is interior,
$r=\nabla F^*(q)$ and $G_Q=F$. Elsewhere $q(r)$ lies in the finite
union of facets. On the measurable set where it belongs to a fixed
facet, its scalar normal component is constant, so its derivative
vanishes almost everywhere there. Indeed, at a differentiability and
density point of a scalar Lipschitz level set, a nonzero derivative
would exclude a cone of fixed positive density from that set. It follows that
\[
 \det D^2G_Q=\boldsymbol1_{\{\nabla F(r)\in Q\}}
                          \det D^2F(r)\quad\text{almost everywhere}.
\]
In our normalization,
$(\ddc F)^n=n!\det D^2F(r)\,dr\,d\tau$.
Changing variables $q=\nabla F(r)$ proves \eqref{eq:fullma}.
If $d\vol_\omega=a_\omega(r,\tau)dr\,d\tau$, then
\[
 d\mu_Q=
 \frac{n!\boldsymbol1_{\{\nabla F(r)\in Q\}}\det D^2F(r)}
      {a_\omega(r,\tau)}d\vol_\omega.
\]
Throughout, $\theta$ remains the Chern curvature.
The density-point assertion applies \cite[Theorem 6.10]{Hunter} to
the locally integrable indicator of each measurable facet level set;
finite global measure is not required.

\section{A rational-map example with nonclosed sampling}
Take $n=2$, $L=\mathcal O(4,4)$, $w=0$, and
\[
 Q=\{(2+t,2-t):-\tfrac12\leq t\leq\tfrac12\}.
\]
Here $b=(2,2)$, $\Lambda=\Z(1,-1)$, $d=1$, and $dq_\Lambda=dt$.
The actual spaces and count are
\[
 V_m=\operatorname{span}\{z_1^{2m+j}z_2^{2m-j}e^m:
                          j\in\Z,\ |j|\leq m/2\},
 \qquad N_m=2\lfloor m/2\rfloor+1.
\]
The limiting mass is one, including odd degrees. The inverse gradient is
\[
 r_1(t)=\log\frac{2+t}{2-t},\qquad r_2(t)=-r_1(t).
\]
The limit averages angular Haar probabilities on
\[
 |z_1|^2=\frac{2+t}{2-t},\qquad
 |z_2|^2=\frac{2-t}{2+t},\qquad -\tfrac12\leq t\leq\tfrac12.
\]
It is carried by a compact three-dimensional portion of $|z_1z_2|=1$.
If $c(p)=p\log(p/4)+(4-p)\log((4-p)/4)$, then
$F^*(2+t,2-t)=2c(2+t)$, giving an explicit envelope.

Adjacent section ratios are $z_1/z_2$. In homogeneous coordinates
$z_j=v_j/u_j$, the map $[v_1u_2:u_1v_2]$ is indeterminate at
$(0,0)$ and $(\infty,\infty)$; it is not a global submersion on $X$.
The norm probabilities concentrate in a transverse radial direction,
as proved above, rather than retaining a fixed fiber probability.

Let $\omega_0$ be the sum of the two Fubini--Study forms and choose
$f=c|z_1|^2/(1+|z_1|^2)$, with $c\ne0$. Then
$\omega=e^f\omega_0$ is positive and
$d\omega=e^fdf\wedge\omega_0\ne0$ where $df\ne0$.
Sampling with $\nu=\omega^2/\int_X\omega^2$ meets the hypotheses and
gives the same mass-one limit. The curvature is still $\theta=4\omega_0$.

\section{An actual multiplier-ideal series}
The preceding theorem prescribes a multiplicative Newton series. We now
compute a distinct, actual multiplier-ideal selection, including its full
class of possibly non-invariant competitors. This calculation uses its own
compact sampling set and norm; it does not apply the preceding whole-space
sampling theorem to different data.

Let $X=\mathbb P^1\times\mathbb P^1$, $L=\mathcal O(1,1)$ with the
product Fubini--Study metric, and $\theta=c_1(L,h)$. In coordinates
$(Z_i:T_i)$ put $z_i=Z_i/T_i$ and
\[
 \psi=\max\{\log|T_1T_2|_h^2,\log|Z_1Z_2|_h^2\},\qquad
 p_+=(0,\infty),\quad p_-=(\infty,0).
\]
These are the only pole points. On the torus let
\[
 x_i=\log|z_i|^2,\qquad s=x_1+x_2,\qquad
 \varphi_X=\log(1+e^{x_1})+\log(1+e^{x_2}).
\]
Then $\psi=\max(0,s)-\varphi_X$. For every real $t>0$,
$\mathcal I(tv)$ consists of germs $f$ for which $|f|^2e^{-tv}$
is locally integrable. Write $v\preceq_{\mathcal I}\psi$ when
$\mathcal I(tv)\subseteq\mathcal I(t\psi)$ for all $t>0$.
The ordinary upper type constraint is $v\leq\psi+C_v$ for some
finite constant. Both constraints concern genuine $\theta$-psh functions.
Every analytic neighborhood used below is connected, and a psh function
invoked as locally integrable is not identically $-\infty$ on that
neighborhood. Restriction to a connected submanifold is either genuine
or identically $-\infty$; this alternative is retained explicitly.

Take the compact sampling set with logarithmic radii in
$D=[-1,2]^2$ and all angles, and use
\[
 w=\frac{x_1^2+x_2^2}{2}-\varphi_X,\qquad
 d\nu=\frac19\,dx_1\,dx_2\,d\tau_1\,d\tau_2.
\]
Each $d\tau_i$ is angular Haar probability. The norm is precisely
$\int_K|\sigma|_{h^m}^2e^{-mw}d\nu$. Define the actual spaces
$V_m=H^0(X,L^m\otimes\mathcal I(m\psi))$ and their $B_m,U_m$
using this norm. The ideal envelope uses $v\leq w$ quasi-everywhere
on $K$, meaning outside a pluripolar set.

\begin{theorem}\label{thm:actual-ideal}
For all genuine $\theta$-psh functions, including non-invariant ones,
\[
 v\preceq_{\mathcal I}\psi\quad\Longleftrightarrow\quad
 v\leq\psi+C_v.
\]
Consequently ideal and ordinary envelopes agree for every common obstacle
problem, in particular for every smooth obstacle on all $X$.
For the stated $K,w$, their envelope is
\[
 P=g(s)-\varphi_X,\qquad
 g(s)=\sup_{0\leq t\leq1}(ts-t^2)
 =\begin{cases}0&s\leq0,\\s^2/4&0\leq s\leq2,\\s-1&s\geq2.\end{cases}
\]
It has its canonical global extension, with $P(p_\pm)=-\infty$.
For every $m\geq1$ the entire selected space is
\begin{equation}\label{eq:actual-three-bands}
 V_m=\operatorname{span}\{Z_1^aT_1^{m-a}Z_2^bT_2^{m-b}:
       0\leq a,b\leq m,\ |a-b|\leq1\},\qquad N_m=3m+1.
\end{equation}
Its Hilbert growth exponent is one. The kernels satisfy the global
lower bound, with extended-real interpretation at the poles,
\begin{equation}\label{eq:actual-lower}
 U_m\geq P+\frac{\log(9m/(2\pi))}{m}-\frac1{4m^2}.
\end{equation}
They converge locally uniformly on $X\setminus\{p_+,p_-\}$,
globally in $L^1$ at rate $O(\log m/m)$,
and extended-pointwise on $X$. Moreover $e^{U_m}\to e^P$
uniformly on $X$. For each fixed $m\geq2$, the additive difference
$U_m-P$ is unbounded above near either pole.

The probability trace measures $N_m^{-1}e^{-mw}B_m\nu$
converge weakly to the probability $\mu$ given by
\[
 \int_X f\,d\mu=\int_0^1\int f(e^{t/2}\tau_1,e^{t/2}\tau_2)
                         \,d\tau_1\,d\tau_2\,dt.
\]
The intrinsic normalization $m^{-1}e^{-mw}B_m\nu$ converges to
$3\mu$, of mass three. Ambient normalization by $m^2$ gives zero.
\end{theorem}

\subsection{All-weight ideals detect the full ordinary constraint}
Near $p_+$ use $u=z_1$, $v=1/z_2$ and the frame $T_1Z_2$.
The two pencil sections become $v,u$, so
\[
 \psi=\log\max(|u|^2,|v|^2)
            -\log(1+|u|^2)-\log(1+|v|^2).
\]
The same formula holds near $p_-$. If a holomorphic germ has order $k$,
its leading homogeneous term and spherical squared average reduce
integrability to $\int_0^\epsilon r^{2k-2t+3}dr$.
The strict threshold is $k+2>t$, with logarithmic divergence at equality.
Thus at either pole
\begin{equation}\label{eq:actual-local-ideal}
 \mathcal I(t\psi)=\mathfrak m^{\max(0,\lfloor t\rfloor-1)};
\end{equation}
away from them it is the full structure sheaf. In particular the
integer degree $m$ imposes order $m-1$.

We use the weighted point extension consequence of Ohsawa--Takegoshi:
for a psh weight $\Phi$ on a bounded ball and a point $a$ in a fixed
smaller ball with $\Phi(a)>-\infty$, there is a holomorphic $f$
with $f(a)=1$ and
\[
 \int |f|^2e^{-\Phi}\,dV\leq C e^{-\Phi(a)}.
\]
The constant is independent of $a$ and $\Phi$, and singular weights
are permitted. Theorem A of \cite{Berndtsson} proves singular-weight hyperplane
extension with the exponent $e^{-\Phi}$ and constant $4\pi$
when the transverse coordinate has modulus less than one.
Use an affine flag through $a$ inside a ball of radius $R$.
Every slice is a connected ball and its weight is finite at $a$.
Scaling its coordinates by $2R$ gives an ambient-to-slice volume
ratio $(2R)^2$. Iterating from the point, whose zero-dimensional
measure is one, gives a valid constant $(16\pi R^2)^n$,
independent of the weight and point. An alternative primary proof
is in \cite[Corollary 12.9, Remark 12.10 and proof of Theorem 12.11]{DemaillyTrieste}.
Here the weight is $\Phi=t\phi$, without an extra factor two.
For a local genuine potential $\phi=v+\varphi_h$, smooth metric
weights do not change multiplier ideals. Shrink the ball inside its
domain so that $\phi$ is bounded above on the closure.

If every germ in $\mathcal I(t\phi)$ has order at least $k$, every
holomorphic function of finite weighted norm on that ball has the same
vanishing jets. Ordinary $L^2$ control from the upper bound for $\phi$
and Cauchy estimates give, on a further smaller ball,
\[
 |f(z)|^2\leq A_{t,k}|z|^{2k}\int |f|^2e^{-t\phi}dV.
\]
Apply point extension at each $z$ with $\phi(z)>-\infty$ to get
$\phi(z)\leq(k/t)\log|z|^2+C_{t,k}$. It holds trivially at
the other points. If $k>0$, a finite $\phi(0)$ would let point
extension produce a function equal to one at the origin, contradicting
its forced vanishing. Thus the origin has the correct value too.
The constants need not be uniform in $t$.

To retain the endpoint coefficient, set
$M(s)=\sup_{|z|\leq e^{s/2}}\phi(z)$. It is finite, nondecreasing
and convex. Indeed the one-variable three-circles maximum on every
complex line is convex, and their supremum is convex; identically
$-\infty$ restrictions do not affect the supremum.
Its limiting left slope $\ell$ is at least every $k(t)/t$ when
the ideal order is $k(t)=\max(0,\lfloor t\rfloor-1)$.
Thus $\ell\geq1$. Convexity gives
\[
 M(s)\leq M(s_0)+\ell(s-s_0)\leq s+C\quad(s<s_0).
\]
It follows that $v\leq\psi+C$ near both poles. On the compact
complement, $\psi$ is finite and continuous and $v$ has a finite
upper bound; a single global constant results. Conversely
$v\leq\psi+C$ immediately implies
$e^{-t\psi}\leq e^{tC}e^{-tv}$ and all the ideal inclusions.
This proves equality of the unrestricted competitor classes before
any obstacle is imposed, and hence equality of their envelopes for
every common obstacle.

\subsection{Pencil rigidity and non-invariant competitors}
Off $p_\pm$ define
\[
 \mathfrak f=[T_1T_2:Z_1Z_2],\qquad
 \psi_{\rm FS}=\log(|T_1T_2|_h^2+|Z_1Z_2|_h^2).
\]
Then $\psi\leq\psi_{\rm FS}\leq\psi+\log2$.
Every admissible $v$ is of the form
\begin{equation}\label{eq:actual-pencil}
 v=\psi_{\rm FS}+u\circ\mathfrak f,\qquad
                   u\in\PSH(\mathbb P^1,\omega_{\rm FS}),
\end{equation}
with canonical extension at the two base points.

To prove this, for $\xi\in\C^*$ the pencil fiber
$C_\xi=\{Z_1Z_2=\xi T_1T_2\}$ is a smooth rational $(1,1)$
curve containing both poles. Its $\theta$-degree is two.
Every genuine restriction of $v$ has a positive curvature measure
of total mass two, and its logarithmic upper bounds force an atom
of mass at least one at each pole. These two atoms exhaust the
measure. The curvature of $\psi|_{C_\xi}$ is exactly these two
unit atoms by Poincare--Lelong. Hence $v-\psi$ is harmonic in
distributions on the compact curve and is constant.
The normalization and logarithmic-atom statement also follow directly
from \cite[III, Sections 2.C and 6]{DemaillyBook}, after replacing
their $\log|z|$ by $\log|z|^2$ and their $\ddc$ by half of it.

In holomorphic torus coordinates $(z_1,\xi=z_1z_2)$, local
integrability and Fubini show that almost every fiber restriction
is genuine. The transverse section $S=\{z_1=1\}$ is a copy of
$\mathbb P^1$; its restriction of $v$ is therefore genuine: almost every good fiber is constant after adding
the metric potential, and hence is finite at its point on $S$.
Put $u=v|_S+\log2$. Since $\theta|_S=\omega_{\rm FS}$ and
$\psi_{\rm FS}|_S=-\log2$, the fiber constants give
\eqref{eq:actual-pencil} almost everywhere on the torus. Both
sides have canonical psh frame potentials, so equality holds
everywhere off the poles. Their logarithmic upper bounds force
value $-\infty$ at the poles. Conversely any such $u$ gives a
positive pulled-back curvature off the poles, and its local
potential is at most $\log(|u_0|^2+|v_0|^2)+C$ there. The
canonical removable psh extension is global and has the required
ordinary upper type. This proves the full description.

On the affine base write $\chi(\xi)=\log(1+|\xi|^2)+u(\xi)$.
For any competitor, $v=\chi(z_1z_2)-\varphi_X$. On the open
interior of $K$ the quasi-everywhere obstacle inequality is an
almost-everywhere inequality. Submean inequalities for the genuine
frame potential and continuity of the obstacle make it pointwise.
Taking $x_1=x_2=s/2$ yields $\chi(\xi)\leq s^2/4$ for
$-2<s<4$, with any argument of $\xi$.

The bound at $s=0$ is $\chi\leq0$ on $|\xi|=1$.
The subharmonic maximum principle gives $\chi\leq0$ on that disk.
At infinity use the genuine local potential
\[
 \chi_\infty(\eta)=\chi(1/\eta)+\log|\eta|^2
                  =\log(1+|\eta|^2)+u(1/\eta).
\]
On $|\eta|=e^{-1}$, the obstacle bound at $s=2$ gives
$\chi_\infty\leq-1$. The same maximum principle on this disk
therefore gives $\chi(\xi)\leq\log|\xi|^2-1$ for $s\geq2$.
Together with the obstacle bound for $0\leq s\leq2$, these are
exactly $\chi(\xi)\leq g(\log|\xi|^2)$ for every $\xi\in\C^*$,
without an angular-invariance hypothesis.
Thus every full competitor is at most $P$. Conversely $g$ is
convex with slopes in $[0,1]$, and defines a base Fubini--Study
psh potential. Also $\psi-1\leq P\leq\psi$ and
$g(s)\leq s^2/4\leq(x_1^2+x_2^2)/2$. Hence $P$ itself is
an admissible global competitor. This identifies the full envelope.

\subsection{The entire selected spaces and exact norms}
At the two poles, the monomial indexed by $(a,b)$ has orders
$a+m-b$ and $b+m-a$. By \eqref{eq:actual-local-ideal} these
must be at least $m-1$, giving exactly $|a-b|\leq1$.
Distinct local exponent pairs cannot cancel; this proves the
entire space \eqref{eq:actual-three-bands}. There are $m+1$
diagonal terms and $m$ in each neighboring band. Ordinary upper
type for each degree-$m$ logarithmic section would instead impose
order $m$ and select only $a=b$. Equality of the all-weight
competitor classes does not identify these finite-degree spaces.

The actual spaces are not multiplicative: the square of an
offset-one section of $V_1$ has offset two and is outside $V_2$.
Their rational maps can recover both $z_1$ and $z_2$ from section
ratios. Their Hilbert growth exponent one is consequently not an
Iitaka-dimension assertion for a graded section ring.

Angular integration makes all monomials orthogonal. For $t\in[0,1]$
set
\[
 I_m(t)=\int_{-1}^2e^{m(tx-x^2/2)}dx
       =e^{mt^2/2}J_m(t),\qquad
 J_m(t)=\int_{-1}^2e^{-m(x-t)^2/2}dx.
\]
The exact original squared norm is
\[
 A_m(a,b)=\frac19 I_m(a/m)I_m(b/m).
\]
Uniformly in $t\in[0,1]$,
\[
 \frac c{\sqrt m}\leq J_m(t)\leq\sqrt{\frac{2\pi}m},\qquad
 J_m(t)=\sqrt{\frac{2\pi}m}
             \bigl(1+O(m^{-1/2}e^{-m/2})\bigr),
\]
where $c=\int_{-1}^1e^{-y^2/2}dy$. The centers are at least
distance one from the endpoints, so each omitted Gaussian tail
is at most $e^{-m/2}/m$. In particular
\begin{equation}\label{eq:actual-norms}
 \frac{c^2}{9m}e^{(a^2+b^2)/(2m)}\leq A_m(a,b)
             \leq\frac{2\pi}{9m}e^{(a^2+b^2)/(2m)}.
\end{equation}

Every basis contribution to $e^{-mw}B_m\nu$ is exactly the
product of the two truncated Gaussian probabilities with centers
$(a/m,b/m)$, times independent angular Haar. They concentrate
uniformly at their centers. Averaging the diagonal and neighboring
bands gives three Riemann sums with the same diagonal limit;
each has asymptotic proportion $1/3$ in $N_m$. This proves the
probability and intrinsic mass-three limits. They are carried
by the contact set $x_1=x_2=t\in[0,1]$, since $g(2t)=t^2$.

\subsection{Global lower bounds and the pole error}
The diagonal sections belong to the actual space. Their upper
norm bound in \eqref{eq:actual-norms} gives
\[
 B_m\geq\frac{9m}{2\pi}e^{-m\varphi_X}
                    \max_{0\leq j\leq m}e^{js-j^2/m}.
\]
The nearest grid point to $s/2$ loses at most $1/(4m^2)$ in
$g(s)$ when $0\leq s\leq2$; outside this interval the
maximizing endpoint is on the grid. This proves
\eqref{eq:actual-lower} globally.

For the upper bound use all three bands. Comparing an offset-one
term to its two neighboring diagonal terms bounds its exponent
by $mg(s)+\min(x_1,1-x_2)$; the other band interchanges the
indices. Set
\[
 R(x)=1+\min(e^{x_1},e^{1-x_2})
             +\min(e^{x_2},e^{1-x_1}).
\]
Counting terms and using the lower norm bound yields
\begin{equation}\label{eq:actual-upper}
 U_m\leq P+\frac{C+2\log(m+1)+\log R(x)}m.
\end{equation}
$R$ extends continuously across coordinate divisors away from
the poles. Near $p_+$ set $A=|u|^2$, $B=|v|^2$.
Its only possibly large term satisfies
$\min(1/B,e/A)\leq(e+1)/(A+B)$; the other is bounded.
The same holds near $p_-$. Thus $R\leq C(1+r^{-2})$ near
either pole, with $r$ a smooth local distance, and $\log R$
is integrable against any smooth volume.

The two bounds prove local uniform convergence off the poles
and the global $L^1$ rate.
For $m\geq2$ every selected section vanishes at both poles,
so $U_m(p_\pm)=-\infty=P(p_\pm)$. The degree-one space is
base-point free and does not affect extended-pointwise convergence.

For a fixed $m\geq2$, the lowest local terms include every
degree-$(m-1)$ binary monomial. Their positive sum of squares
is comparable to $r^{2(m-1)}$ in all directions. Therefore
\[
 U_m=(1-1/m)\log r^2+O_m(1),\qquad
 P=\log r^2+O(1).
\]
The additive difference diverges to $+\infty$ on punctured
approaches to either pole. Finally \eqref{eq:actual-upper},
$P\leq\log r^2+C$ and $R\leq C/r^2$ give
$e^{U_m}\leq C r^{2-2/m}\leq Cr$ for $m\geq2$, uniformly
in $m$ and $r<1$. Also $e^P\leq Cr^2$. Together with uniform
convergence on compacta off the poles these bounds prove the
asserted uniform exponential convergence.

\section{Scope and excluded cases}
Without the integral offset a rational affine hull can lack lattice
points in infinitely many degrees: a half-integral singleton is empty
in every odd degree. Our hypotheses allow finitely many empty initial
spaces but guarantee nonzero spaces in every sufficiently large degree.
Every nonempty allowed polytope has positive relative volume, including
a singleton with the stated convention.

For the Newton theorem, polytopes touching the degree-box boundary,
nonconvex logarithmic
weights, non-invariant sampling or obstacles, singular sampling
probabilities and arbitrary compact sampling sets are not covered.
The actual ideal example uses the specific compact sampling set, weight
and singularity computed above. Its all-weight competitor classes agree,
while its finite-degree ideal spaces have neighboring bands and are not
multiplicative. The two independent theorems completely prove these
explicit sufficient-condition answers to P20. Arbitrary selecting or
multiplier-ideal series, universal ideal-envelope identities and the
broader original programme remain open.
\begin{thebibliography}{9}
\bibitem{BEGZ} S. Boucksom, P. Eyssidieux, V. Guedj and A. Zeriahi,
\emph{Monge--Amp\`ere equations in big cohomology classes},
\href{https://arxiv.org/abs/0812.3674}{arXiv:0812.3674}, Section 1.2.
\bibitem{Kinnunen} J. Kinnunen, \emph{Sobolev Spaces},
\href{https://math.aalto.fi/~jkkinnun/files/sobolev_spaces.pdf}{course notes},
Section 3.4, pp.~78--81.
\bibitem{Hunter} J. K. Hunter, \emph{Differentiation},
\href{https://www.math.ucdavis.edu/~hunter/m206/ch5_measure_notes.pdf}{measure-theory notes},
Theorem 6.10, pp.~68--70.

\bibitem{Berndtsson} B. Berndtsson,
\emph{The extension theorem of Ohsawa--Takegoshi and the theorem of
Donnelly--Fefferman}, Ann. Inst. Fourier 46 (1996), 1083--1094,
\href{https://www.numdam.org/item/10.5802/aif.1541.pdf}{primary text},
Theorem A and Section 2.
\bibitem{DemaillyTrieste} J.-P. Demailly,
\emph{Multiplier ideal sheaves and analytic methods in algebraic geometry},
Trieste lecture notes (2000),
\href{https://www-fourier.univ-grenoble-alpes.fr/~demailly/manuscripts/trieste.pdf}{primary text},
Section 12, especially Corollary 12.9, Remark 12.10 and Theorem 12.11.
\bibitem{DemaillyBook} J.-P. Demailly,
\emph{Complex Analytic and Differential Geometry},
\href{https://people.math.harvard.edu/~demarco/Math274/Demailly_ComplexAnalyticDiffGeom.pdf}{primary text},
I.4--5 and III.2.C, III.6.
\end{thebibliography}
\end{document}
